\xw{defer to some place later or appendix or discussion.}
For each Pauli component, knowing the sign of one link parameter
would resolve the ambiguity, since the other two signs follow from
the measured products. In particular, a known reference link
\(i_0\) satisfying \(A_{i_0}>1/2\) would imply
\(q_{X,\ell _0},q_{Y,\ell _0},q_{Z,\ell _0}>0\), allowing all three channels
to be identified from the three sets of pairwise products.
In this work, we do not assume such prior knowledge. 

\section{Proof of Lemma~\ref{lem:double-selection-statistics}}
\label{app:double-selection}

We label the Bell states by two bits \(a,b\in\{0,1\}\):
\begin{equation}
\ket{\beta_{ab}}
:=
(I\otimes X^aZ^b)\ket{\Phi^+}.
\end{equation}
The labels \((0,0),(1,0),(1,1),(0,1)\) correspond to
\(\ket{\Phi^+},\ket{\Psi^+},\ket{\Psi^-},\ket{\Phi^-}\),
respectively.

Let \((a_i,b_i)\) denote the Bell label of the \(i\)th input pair.
After the bilateral CNOTs from pair \(1\) to pair \(2\) and from
pair \(3\) to pair \(2\), the output labels are
\begin{equation}
\begin{aligned}
(a_1',b_1')
&=(a_1,\;b_1\oplus b_2),\\
(a_2',b_2')
&=(a_1\oplus a_2\oplus a_3,\;b_2),\\
(a_3',b_3')
&=(a_3,\;b_2\oplus b_3).
\end{aligned}
\label{eq:double-selection-bell-labels}
\end{equation}
Thus, coincidence in the \(Z\) measurements occurs if and only if
\(a_1\oplus a_2\oplus a_3=0\), while coincidence in the \(X\)
measurements occurs if and only if \(b_2\oplus b_3=0\).

Define
\begin{equation}
\begin{aligned}
x&:=\mathbb E[(-1)^b]=2(q_1+q_2)-1,\\
y&:=\mathbb E[(-1)^{a+b}]=2(q_1+q_3)-1,\\
z&:=\mathbb E[(-1)^a]=2(q_1+q_4)-1.
\end{aligned}
\label{eq:double-selection-correlations}
\end{equation}
Let
\[
U:=(-1)^{a_1+a_2+a_3},
\qquad
V:=(-1)^{b_2+b_3}.
\]
By independence of the three input pairs,
\begin{equation}
\mathbb E[U]=z^3,
\qquad
\mathbb E[V]=x^2,
\qquad
\mathbb E[UV]=zy^2.
\end{equation}
The coincidence probabilities therefore satisfy
\begin{equation}
\begin{aligned}
p_Z
&:=
\Pr(z_A=z_B)
=
\mathbb E\left[\frac{1+U}{2}\right]
=
\frac{1+z^3}{2},\\
p_X
&:=
\Pr(x_A=x_B)
=
\mathbb E\left[\frac{1+V}{2}\right]
=
\frac{1+x^2}{2},\\
p_{XZ}
&:=
\Pr(z_A=z_B,\;x_A=x_B)\\
&=
\mathbb E\left[\frac{(1+U)(1+V)}{4}\right]
=
\frac{1+z^3+x^2+zy^2}{4}.
\end{aligned}
\label{eq:double-selection-probabilities}
\end{equation}

Since \(q_1>1/2\), all three quantities \(x,y,z\) are strictly
positive. Hence, these relations can be uniquely inverted:
\begin{equation}
\begin{aligned}
x&=\sqrt{2p_X-1},\\
z&=\sqrt[3]{2p_Z-1},\\
y&=
\sqrt{
\frac{4p_{XZ}-2p_X-2p_Z+1}{z}
}.
\end{aligned}
\label{eq:double-selection-inversion}
\end{equation}
The positive square-root branches are fixed by the assumption
\(q_1>1/2\), which also ensures \(z>0\).
The Bell-state probabilities are then recovered as
\begin{equation}
\begin{bmatrix}
q_1\\q_2\\q_3\\q_4
\end{bmatrix}
=
\frac{1}{4}
\begin{bmatrix}
1+x+y+z\\
1+x-y-z\\
1-x+y-z\\
1-x-y+z
\end{bmatrix}.
\label{eq:double-selection-state-reconstruction}
\end{equation}
This proves that the three coincidence probabilities uniquely
determine the input Bell-diagonal state.

It remains to determine the conditional output fidelity.
By Eq.~\eqref{eq:double-selection-bell-labels}, the first output
pair is in \(\ket{\Phi^+}\) when
\[
a_1=0,
\qquad
b_1=b_2.
\]
Combining these conditions with the acceptance conditions gives
\[
a_2=a_3,
\qquad
b_1=b_2=b_3.
\]
Therefore, the unnormalized weight of \(\ket{\Phi^+}\) in the
accepted output is
\begin{equation}
\begin{aligned}
N_{\Phi^+}
&=
q_1(q_1^2+q_2^2)+q_4(q_4^2+q_3^2)\\
&=
q_1^3+q_4^3+q_1q_2^2+q_4q_3^2.
\end{aligned}
\end{equation}
Since \(p_{XZ}\geq q_1^3>0\), normalization yields
\begin{equation}
\begin{aligned}
F_{\mathrm{out}}
&=
\frac{N_{\Phi^+}}{p_{XZ}}\\
&=
\frac{
4\left(q_1^3+q_4^3+q_1q_2^2+q_4q_3^2\right)
}{
1+z^3+x^2+zy^2
},
\end{aligned}
\end{equation}
which establishes
Eq.~\eqref{eq:double-selection-output-fidelity}.
\section{Details for Modified \mergecast in Figure \ref{fig:modified-mergecast}}
Below, we present the quantum states at slices \(1\)--\(4\) of
Fig.~\ref{fig:modified-mergecast}. Let \(R\) denote the qubit retained
by \(A_1\), \(T\) the transmitted half of the Bell pair, and \(A\) the
auxiliary qubit sent by \(A_2\). We use the tensor-product ordering
\((R,T,A)\). Since the corresponding three-qubit Pauli--Liouville
vectors contain \(64\) entries, we list only the potentially nonzero
components with respect to the ordered sparse basis
\begin{equation}
\mathcal B_3:=
\left(
III,IIZ,XXI,XXZ,YYI,YYZ,ZZI,ZZZ
\right),
\end{equation}
where, for example, \(XXZ:=X_R\otimes X_T\otimes Z_A\).

At slice \(1\), the joint state is
\begin{equation}
\rho_{RTA}^{(1)}
=
\ket{\Phi^+}\!\bra{\Phi^+}_{RT}
\otimes
\ket{0}\!\bra{0}_A,
\end{equation}
whose nonzero Pauli--Liouville components are
\begin{equation}
\left.\left|\rho_{RTA}^{(1)}\right\rangle\right\rangle_{\mathcal B_3}
=
\begin{bmatrix}
1\\
1\\
1\\
1\\
-1\\
-1\\
1\\
1
\end{bmatrix}.
\end{equation}

After \(T\) and \(A\) pass through \(\mathcal P_1\) and
\(\mathcal P_2\), respectively, the state at slice \(2\) is
\begin{equation}
\rho_{RTA}^{(2)}
=
\left(
\mathcal I_R\otimes\mathcal P_1\otimes\mathcal P_2
\right)
\left(\rho_{RTA}^{(1)}\right),
\end{equation}
with
\begin{equation}
\left.\left|\rho_{RTA}^{(2)}\right\rangle\right\rangle_{\mathcal B_3}
=
\begin{bmatrix}
1\\
q_{Z2}\\
q_{X1}\\
q_{X1}q_{Z2}\\
-q_{Y1}\\
-q_{Y1}q_{Z2}\\
q_{Z1}\\
q_{Z1}q_{Z2}
\end{bmatrix}.
\end{equation}

Node \(C\) then applies
\(U_{\mathrm C}:=\operatorname{CNOT}_{A\rightarrow T}\).
Immediately after the CNOT, the state is
\begin{equation}
\widetilde{\rho}_{RTA}^{(3)}
=
U_{\mathrm C}\rho_{RTA}^{(2)}U_{\mathrm C}^{\dagger},
\end{equation}
and its nonzero components are
\begin{equation}
\left.\left|\widetilde{\rho}_{RTA}^{(3)}
\right\rangle\right\rangle_{\mathcal B_3}
=
\begin{bmatrix}
1\\
q_{Z2}\\
q_{X1}\\
q_{X1}q_{Z2}\\
-q_{Y1}q_{Z2}\\
-q_{Y1}\\
q_{Z1}q_{Z2}\\
q_{Z1}
\end{bmatrix}.
\end{equation}
Here, the CNOT exchanges the Pauli components
\[
YYI\longleftrightarrow YYZ,
\qquad
ZZI\longleftrightarrow ZZZ,
\]
while leaving the other components in \(\mathcal B_3\) unchanged.

Tracing out the auxiliary qubit removes all Pauli strings whose
auxiliary component is \(Z\). Therefore, at slice \(3\),
\begin{equation}
\rho_{RT}^{(3)}
:=
\operatorname{Tr}_A
\left(\widetilde{\rho}_{RTA}^{(3)}\right)
\end{equation}
has the restricted Pauli--Liouville representation
\begin{equation}
\left.\left|\rho_{RT}^{(3)}\right\rangle\right\rangle
=
\begin{bmatrix}
1\\
q_{X1}\\
-q_{Y1}q_{Z2}\\
q_{Z1}q_{Z2}
\end{bmatrix},
\end{equation}
with respect to the ordered basis
\((II,XX,YY,ZZ)\).

Finally, after \(T\) passes through \(\mathcal P_3\), the state shared
between \(A_1\) and \(A_3\) at slice \(4\) is
\begin{equation}
\rho_{13}^{(0)}
=
\left(\mathcal I_R\otimes\mathcal P_3\right)
\left(\rho_{RT}^{(3)}\right),
\end{equation}
and hence
\begin{equation}
\left.\left|\rho_{13}^{(0)}\right\rangle\right\rangle
=
\begin{bmatrix}
1\\
q_{X1}q_{X3}\\
-q_{Y1}q_{Z2}q_{Y3}\\
q_{Z1}q_{Z2}q_{Z3}
\end{bmatrix}.
\label{eq:mergecast-slice-state}
\end{equation}